\newif\ifglobal

%%%%%%%%%%%%%%%%%%%%%%%
%%% Compile to view %%%
%%%%%%%%%%%%%%%%%%%%%%%

%%%%%%%%%%%%%%%%%%%%%%%%%%%%%%%%%%%%%%%%%%%%%%%%%%%
%%% Readme:                                     %%%
%%% https://www.overleaf.com/read/bwdjvfjksvjy/ %%%
%%%%%%%%%%%%%%%%%%%%%%%%%%%%%%%%%%%%%%%%%%%%%%%%%%%

% >>> M?: marks a module setting number or important notice

% >>> M4: Set {./relative-path} to external files for this module <<<
\def \modulefiles {./08-INTMTRX}
% To add an external file, use:
% (Replace '---' with actual filename)
% '\input{\modulefiles/tables/---.tex}' for tables
% '\includegraphics{\modulefiles/figures/---}' for figures
% '\subfile{\modulefiles/---}' for register files

% >>> M1: This module is in global project? <<<
% 'YES' - uncomment; 'NO' - comment out
\globaltrue

\ifglobal
    \documentclass[main\_\_EN.tex]{subfiles}
   
\else
    \documentclass[openany, 10pt]{book}

    % >>> M2: In ./00-shared/preamble-trm-module, also find
    %         settings for visibility of todo notes, watermark <<<
    \usepackage{preamble-trm-module}

    % >>> M3: Temporary labels in English,
    %         you can add your own labels there <<<
    \myexternaldocument{./00-shared/temp-labels-en}

\fi %\ifglobal


\setcounter{secnumdepth}{4}


% >>> M5: If this module needs any updates in
% './00-shared/preamble-trm-module.sty'
% (include additional package, etc.), contact your module PM
% or create the file './00-shared/changelog.tex' make the
% necessary updates yourself and document them in this file <<<

% Make text left-aligned and allow latex to break pages naturally, comment out for full justification
\raggedright
\raggedbottom


\begin{document}


% Sets language into which pre-defined bits of text are translated
\selectlanguage{english}
\setenglishformatting

% Do not delete this block
\ifglobal
% Add the following front matter if
% this module is outside of global project
\else
    \InputIfFileExists{readme}{}{}
    \listoftodos
    {\let\clearpage\relax \listofdonetodos}
    \tableofcontents
    %\listoftables
    %\listoffigures
\fi



%%%%%%%%%%%%%%%%%%%%%%%%%
%%% TRM Module Begins %%%
%%%%%%%%%%%%%%%%%%%%%%%%%

\hypertarget{intmtrx}{}
\chapter{Interrupt Matrix}
\label{mod:intmtrx}


\section{Overview}

The interrupt matrix embedded in ESP32-S2 independently allocates peripheral interrupt sources to the CPU peripheral interrupts, so as to timely inform the CPU to process the interrupts once the interrupt signals are generated. This flexible function is applicable to a variety of application scenarios.

\section{Features}

\begin{itemize}
    \item Accept 95 peripheral interrupt sources as input
    \item Generate 26 peripheral interrupts to the CPU as output
    \item Disable CPU non-maskable interrupt (NMI) sources
    \item Query current interrupt status of peripheral interrupt sources
\end{itemize}

The structure of the interrupt matrix is shown in Figure \ref{figure:InterruptMatrixStructure}.

\begin{figure}[h! ]
    \centering
    \fbox{\includegraphics[width=0.9\textwidth]{\modulefiles/figures/InterruptMatrixStructure.eps}}
    \caption{Interrupt Matrix Structure}
    \label{figure:InterruptMatrixStructure}
\end{figure}

\section{Functional Description}

\subsection{Peripheral Interrupt Sources}

ESP32-S2 has 95 peripheral interrupt sources in total, all of which can be allocated to the CPU. For the peripheral interrupt sources and their configuration/status registers, please refer to Table \ref{table:PeripheralInterruptConfigStatusSource}.

\clearpage
\begin{landscape}
\subfile{\modulefiles/PeripheralInterruptConfigStatusSource_en.tex}

\end{landscape}
\restoregeometry

\subsection{CPU Interrupts}

The CPU has 32 interrupts, including 26 peripheral interrupts and six internal interrupts. Table \ref{table:CPUInterrupt} lists all the interrupts.

\begin{itemize}
    \item Peripheral Interrupts:
    \begin{itemize}
        \item Level-triggered interrupts: triggered by high level signal. The interrupt sources should hold the level till the CPU handles the interrupts.
        \item Edge-triggered interrupts: triggered on rising edge. The CPU responds to this kind of interrupts immediately.
        \item NMI interrupt: once triggered, the NMI interrupt can not be masked by software using the CPU registers.
    \end{itemize}
    \item Internal Interrupts:
        \begin{itemize}
        \item Timer interrupts: triggered by internal timers and are used to generate periodic interrupts.
        \item Software interrupts: triggered when software writes to special registers.
        \item Profiling interrupt: triggered for performance monitoring and analysis.
        \end{itemize}

\end{itemize}

ESP32-S2 supports the above-mentioned 32 interrupts at six levels as shown in the table below. A higher level corresponds to a higher priority. NMI has the highest interrupt priority and once triggered, the CPU must handle such interrupt.

\subfile{\modulefiles/CPUInterrupt_en.tex}


\subsection{Allocate Peripheral Interrupt Source to CPU Interrupt}

In this section, the following terms are used to describe the operation of the interrupt matrix.

\begin{itemize}

\item Source\_X: stands for a particular peripheral interrupt source, wherein, X means the number of this interrupt source in Table \ref{table:PeripheralInterruptConfigStatusSource}.

\item INTERRUPT\_PRO\_X\_MAP\_REG: stands for a peripheral interrupt configuration register, corresponding to the peripheral interrupt source Source\_X. In Table \ref{table:PeripheralInterruptConfigStatusSource}, the registers listed in column "Configuration Register” correspond to the peripheral interrupt sources listed in column “Source”. For example, the configuration register for source UHCI0\_INTR is \hyperref[regdesc:INTERRUPTPROUHCI0INTRMAPREG]{INTERRUPT\_PRO\_UHCI0\_INTR\_MAP\_REG}.

\item Interrupt\_P: stands for the CPU peripheral interrupt numbered as Num\_P. The value of Num\_P can be 0 \verb+~+ 5, 8 \verb+~+ 10, 12 \verb+~+ 14, 17 \verb+~+ 28 and 30 \verb+~+ 31 (see Table \ref{table:CPUInterrupt}).

\item Interrupt\_I: stands for the CPU internal interrupt numbered as Num\_I. The value of Num\_I can be 6, 7, 11, 15, 16 and 29 (see Table \ref{table:CPUInterrupt}).

\end{itemize}

\subsubsection{Allocate one peripheral interrupt source Source\_X to CPU}

Setting the corresponding configuration register INTERRUPT\_PRO\_X\_MAP\_REG of Source\_X to Num\_P will allocate this interrupt source to Interrupt\_P. Num\_P here can be any value from 0 \verb+~+ 5, 8 \verb+~+ 10, 12 \verb+~+ 14, 17 \verb+~+ 28 and 30 \verb+~+ 31. Note that one CPU interrupt can be shared by multiple peripherals.

\subsubsection{Allocate multiple peripheral interrupt sources Source\_X\regindex{n} to CPU}

Setting the corresponding configuration register INTERRUPT\_PRO\_X\regindex{n}\_MAP\_REG of each interrupt source to the same Num\_P will allocate all the sources to the same Interrupt\_P. Any of these sources will trigger CPU Interrupt\_P. When an interrupt signal is generated, software should check the interrupt status registers to figure out which peripheral the signal comes from.

\subsubsection{Disable CPU peripheral interrupt source Source\_X}
Setting the corresponding configuration register INTERRUPT\_PRO\_X\_MAP\_REG of the source to any Num\_I will disable this interrupt Source\_X. The choice of Num\_I does not matter, as none of Num\_I is connected to the CPU. Therefore this functionality can be used to disable peripheral interrupt sources.


\subsection{Disable CPU NMI Interrupt Sources}

The interrupt matrix is able to mask all peripheral interrupt sources allocated to CPU No.14 NMI interrupt using hardware, depending on the internal signal \hyperref[fielddesc:INTERRUPTPRONMIMASKHW]{INTERRUPT\_PRO\_NMI\_MASK\_HW}. The signal comes from "Interrupt Reg" register configuration submodule inside interrupt matrix, see Figure \ref{figure:InterruptMatrixStructure}. If the signal is set to high level, CPU will not respond to NMI interrupt.


\subsection{Query Current Interrupt Status of Peripheral Interrupt Source}

Current interrupt status of a peripheral interrupt source can be read via the bit value in \hyperref[regdesc:INTERRUPTPROINTRSTATUSREG0REG]{INTERRUPT\_PRO\_INTR\_STATUS\_REG\_\regindex{n}}. For the mapping between \hyperref[regdesc:INTERRUPTPROINTRSTATUSREG0REG]{INTERRUPT\_PRO\_INTR\_STATUS\_REG\_\regindex{n}} and peripheral interrupt sources, please refer to Table \ref{table:PeripheralInterruptConfigStatusSource}.


\section{Base Address}\label{subsection:interrupt-matrix-base-address}

Users can access interrupt matrix with the base address, which can be seen in the following table. For more information about accessing peripherals from different buses please see Chapter \ref{mod:sysmem}: \textit{\nameref{mod:sysmem}}.


\begin{longtable}[c]{ | l | C{3cm} | }
\caption{Interrupt Matrix Base Address} \label{table:interrupt-matrix-base-address}
\\ \hline
\rowcolor{lightgray}
Bus to Access Peripheral & Base Address \\
    \hline
PeriBUS1 & \hexprefix{}3F4C2000 \\
    \hline
\end{longtable}

\hypertarget{intmtrx-reg-summ}{}
\section{Register Summary}

The address in the following table represents the address offset (relative address) with the respect to the peripheral base address, not the absolute address. For detailed information about the interrupt matrix base address, please refer to Section \ref{subsection:interrupt-matrix-base-address}.

\clearpage
\begin{landscape}
\subfile{\modulefiles/interrupt_summary_en.tex}

\end{landscape}
\restoregeometry

\section{Registers}
The address in the following part represents the address offset (relative address) with respect to the peripheral base address, not the absolute address. For detailed information about the interrupt matrix base address, please refer to Section \ref{subsection:interrupt-matrix-base-address}.
\subfile{\modulefiles/interrupt_reg_en.tex}


\end{document}
